
\section{Isotropic Heisenberg model and DFT mapping}

We describe N localized spin-$1/2$ moments using the isotropic Heisenberg Hamiltonian
\begin{equation}
    \hat H =
    -\sum_{i<j}^{N} J_{ij}\,
    \hat{\mathbf S}_i\cdot\hat{\mathbf S}_j ,
    \label{eq:heisenberg_general}
\end{equation}
where all pairs are allowed to interact. With this convention,
$J_{ij}>0$ denotes ferromagnetic (FM) coupling and $J_{ij}<0$
antiferromagnetic (AFM) coupling. The spin operators are dimensionless
($\hbar=1$).




\subsection{Two spins: diradical}

For two spin-$1/2$ centers,
\begin{equation}
    \hat H=-J\,\hat{\mathbf S}_1\cdot\hat{\mathbf S}_2 .
\end{equation}
Using
\begin{equation}
    \hat{\mathbf S}_1\cdot\hat{\mathbf S}_2
    =
    \hat S_{1z}\hat S_{2z}
    +
    \frac{1}{2}
    \left(
        \hat S_{1+}\hat S_{2-}
        +
        \hat S_{1-}\hat S_{2+}
    \right),
\end{equation}
the Hamiltonian in the basis
$\{
|\uparrow\uparrow\rangle,
|\uparrow\downarrow\rangle,
|\downarrow\uparrow\rangle,
|\downarrow\downarrow\rangle
\}$
is
\begin{equation}
    H=
    -J
    \begin{pmatrix}
        1/4 & 0 & 0 & 0\\
        0 & -1/4 & 1/2 & 0\\
        0 & 1/2 & -1/4 & 0\\
        0 & 0 & 0 & 1/4
    \end{pmatrix}.
\end{equation}

\noindent The eigenstates are the triplet

\begin{equation}
    |T_+\rangle &= |\uparrow\uparrow\rangle,\quad |T_0\rangle &=
    \frac{1}{\sqrt{2}}
    \left(
        |\uparrow\downarrow\rangle+
        |\downarrow\uparrow\rangle
    \right),\quad |T_-\rangle &=|\downarrow\downarrow\rangle,
    \qquad
    E_T=-\frac{J}{4}.
\end{equation}

and the singlet
\begin{equation}
    |S\rangle=
    \frac{1}{\sqrt{2}}
    \left(
        |\uparrow\downarrow\rangle-
        |\downarrow\uparrow\rangle
    \right),
    \qquad
    E_S=\frac{3J}{4}.
\end{equation}
Thus,
\begin{equation}
    E_S-E_T=J.
\end{equation}

In a collinear broken-symmetry DFT calculation,
$|\uparrow\uparrow\rangle$ is a pure triplet state, whereas
$|\uparrow\downarrow\rangle$ is not a pure singlet:
\begin{equation}
    |\uparrow\downarrow\rangle
    =
    \frac{1}{\sqrt{2}}
    \left(
        |T_0\rangle+|S\rangle
    \right).
\end{equation}
Within the ideal spin-$1/2$ mapping,
\begin{align}
    E_{\uparrow\uparrow} &= E_0-\frac{J}{4},\\
    E_{\uparrow\downarrow} &= E_0+\frac{J}{4},
\end{align}
and therefore
\begin{equation}
    \boxed{
    J=2\left(
        E_{\uparrow\downarrow}
        -E_{\uparrow\uparrow}
    \right)}.
    \label{eq:J_two_spin_BS}
\end{equation}

\noindent More generally, spin contamination of the broken-symmetry state can be
included through the Yamaguchi projection. In the present convention,
\begin{equation}
    \boxed{
    J=
    -2\,
    \frac{E_{\rm HS}-E_{\rm BS}}
    {
    \langle \hat S^2\rangle_{\rm HS}
    -
    \langle \hat S^2\rangle_{\rm BS}
    }}.
    \label{eq:yamaguchi_our_convention}
\end{equation}


\subsection{Four spins: molecular dimer}

For four spin-$1/2$ centers with all-to-all connectivity,
\begin{equation}
    \hat H =
      -J_{12}\hat{\mathbf S}_1\cdot\hat{\mathbf S}_2
      -J_{13}\hat{\mathbf S}_1\cdot\hat{\mathbf S}_3
      -J_{14}\hat{\mathbf S}_1\cdot\hat{\mathbf S}_4
      -J_{23}\hat{\mathbf S}_2\cdot\hat{\mathbf S}_3
      -J_{24}\hat{\mathbf S}_2\cdot\hat{\mathbf S}_4
      -J_{34}\hat{\mathbf S}_3\cdot\hat{\mathbf S}_4.
    \label{eq:four_spin_H}
\end{equation}
\noindent The spectrum consists of
\begin{center}
\begin{tabular}{ccc}
\hline
Total spin & Multiplicity & Number of multiplets\\
\hline
$S=0$ & $1$ & $2$\\
$S=1$ & $3$ & $3$\\
$S=2$ & $5$ & $1$\\
\hline
\end{tabular}
\end{center}
corresponding to two singlets, three triplets, and one quintet.
For arbitrary $J_{ij}$, states with the same total spin can have
different energies; the eigenenergies are therefore obtained by
diagonalizing Eq.~\eqref{eq:four_spin_H}. The quintet energy is
\begin{equation}
    E_{S=2}
    =
    -\frac14
    \left(
    J_{12}+J_{13}+J_{14}+J_{23}+J_{24}+J_{34}
    \right),
\end{equation}
while the two singlet and three triplet energies depend on the complete
set of six exchange constants.

The exchange constants are obtained by mapping the energies of collinear broken-symmetry DFT configurations onto the Heisenberg model. For a collinear configuration, each local spin is assigned a projection
\begin{equation}
    \sigma_i=
    \begin{cases}
        +1,& \uparrow,\\
        -1,& \downarrow.
    \end{cases}
\end{equation}
Using
\begin{equation}
    \hat{\mathbf S}_i\cdot\hat{\mathbf S}_j
    =
    \hat S_{iz}\hat S_{jz}
    +
    \frac12
    \left(
        \hat S_i^+\hat S_j^-
        +
        \hat S_i^-\hat S_j^+
    \right),
\end{equation}
the spin-flip terms generate a different, orthogonal collinear configuration and therefore do not contribute to the diagonal expectation value. Hence,
\begin{equation}
    \left\langle
    \hat{\mathbf S}_i\cdot\hat{\mathbf S}_j
    \right\rangle_{\rm collinear}
    =
    \left\langle
    \hat S_{iz}\hat S_{jz}
    \right\rangle
    =
    \frac14 \sigma_i \sigma_j
    =
    \begin{cases}
        +\frac14, & \text{parallel spins},\\[2pt]
        -\frac14, & \text{antiparallel spins}.
    \end{cases}
\end{equation}
The energy associated with each DFT spin configuration $k$ is therefore mapped as
\begin{equation}
    \boxed{
    E_k =
    E_0
    -
    \frac14\sum_{i<j}
    J_{ij}\,
    \sigma_i^{(k)}\sigma_j^{(k)}
    }.
    \label{eq:DFT_four_spin_mapping}
\end{equation}

Importantly, Eq.~\eqref{eq:DFT_four_spin_mapping} gives the diagonal expectation value of the Heisenberg Hamiltonian for a given collinear configuration; it is not an eigenvalue equation. The off-diagonal terms proportional to $\hat S_i^+\hat S_j^-$ and $\hat S_i^-\hat S_j^+$ are retained when the full Heisenberg Hamiltonian is constructed and diagonalized.